%Polytechnic Institute of Viseu
%School of Technology and Management of Viseu
%File Template Dissertation Created by Gilberto Rouxinol on 2023
%https://www.overleaf.com/
%Copyright © https://ctan.org
%Copyright © MIT License 2023 Gilberto Rouxinol

% --- ESTE FICHEIRO É CHAMADO PELO main.tex ---
% --- NÃO PODE CONTER \documentclass ---

%Texto
\usepackage[utf8]{inputenc}
\usepackage[portuguese]{babel}

%Texto cor
\usepackage{xcolor}

%Bibliografia
\usepackage[backend=biber, style=numeric, sorting=ynt]{biblatex}
\addbibresource{G_RefBiblio.bib}
\usepackage{csquotes}

%Configuração da página
\usepackage[a4paper,left=3.0cm,right=3.0cm,top=2.5cm,bottom=2.5cm]{geometry}

%Alinhar equações
\usepackage{amsmath}

%Fontes ou símbolos matemáticos
\usepackage{amssymb}

%Referências de figuras
\usepackage{graphicx}
\usepackage[export]{adjustbox}

%Colocar subfiguras
\usepackage{caption}
\usepackage{subcaption}

%Escrever unidades do SI
\usepackage{siunitx}

%Acrescentar pdfs
\usepackage{pdfpages}

%Fixar figuras
\usepackage{float}

% --- PACOTE TIKZ PARA DESENHOS VETORIAIS ---
\usepackage{tikz}
\usetikzlibrary{arrows.meta, positioning, calc, patterns}

% --- INÍCIO: PACOTE E ESTILOS DE CÓDIGO (Listings) ---
\usepackage{listings} % Pacote principal para o código

% --- Definir as cores ---
\definecolor{codegreen}{rgb}{0,0.6,0}
\definecolor{codegray}{rgb}{0.5,0.5,0.5}
\definecolor{codepurple}{rgb}{0.58,0,0.82}
\definecolor{backcolour}{rgb}{0.95,0.95,0.92}

% --- Configuração global do listings (IMPORTANTE para acentos e underscores) ---
\lstset{
    literate={á}{{\'a}}1 {ã}{{\~a}}1 {ç}{{\c{c}}}1 {é}{{\'e}}1 {ê}{{\^e}}1 {í}{{\'i}}1 {ó}{{\'o}}1 {õ}{{\~o}}1 {ú}{{\'u}}1 {Á}{{\'A}}1 {Ã}{{\~A}}1 {Ç}{{\c{C}}}1 {É}{{\'E}}1 {Ê}{{\^E}}1 {Í}{{\'I}}1 {Ó}{{\'O}}1 {Õ}{{\~O}}1 {Ú}{{\'U}}1
    {Ø}{{\O}}1 {²}{{$^2$}}1 {←}{{$\leftarrow$}}1 {λ}{{$\lambda$}}1 {√}{{$\sqrt{}$}}2 {β}{{$\beta$}}1 {φ}{{$\phi$}}1 {π}{{$\pi$}}1 {μ}{{$\mu$}}1
    {₀}{{$_0$}}1 {₁}{{$_1$}}1 {₂}{{$_2$}}1
    {ª}{{$^a$}}1 {º}{{$^o$}}1
    {_}{\textunderscore}1 % Trata o underscore como texto normal
}

% --- Estilo para Código Python ---
\lstdefinestyle{PythonStyle}{
    language=Python,
    backgroundcolor=\color{backcolour},   
    commentstyle=\color{codegreen},
    keywordstyle=\color{magenta},
    numberstyle=\tiny\color{codegray},
    stringstyle=\color{codepurple},
    basicstyle=\ttfamily\footnotesize,
    breakatwhitespace=false,         
    breaklines=true,                 
    captionpos=b,                    
    keepspaces=true,                 
    numbers=left,                    
    numbersep=5pt,                  
    showspaces=false,                
    showstringspaces=false,
    showtabs=false,                  
    tabsize=2,
    morekeywords={self, Falso, Verdadeiro, None} % Keywords extra
}

% --- Estilo para Pseudocódigo ---
\lstdefinestyle{PseudoStyle}{
    language=Bash, % Um truque para highlighting básico (trata # como comentário)
    backgroundcolor=\color{backcolour},   
    commentstyle=\color{codegreen},
    keywordstyle=\color{blue}, % Keywords em azul
    stringstyle=\color{codepurple},
    basicstyle=\ttfamily\footnotesize,
    breakatwhitespace=true,         
    breaklines=true,                 
    captionpos=b,                    
    keepspaces=true,                 
    numbers=none, % Sem números de linha para pseudocódigo
    showspaces=false,                
    showstringspaces=false,
    showtabs=false,                  
    tabsize=2,
    % Lista de keywords para o seu pseudocódigo
    morekeywords={ALGORITMO, DADOS, ENTRADA, PARAMETROS, SE, ENTÃO, SENÃO, FIM, PARA, ATÉ, SAIR, DO, CICLO, FUNCAO, RETORNAR, VERIFICAR_DUCTILIDADE, MAX, Otimizar_Armadura_Viga, Dimensionar_Viga_com_Recalculo, Calcular_Resistencias, Calcular_Momento_Reduzido, Calcular_Area_Aco, Extrair_Diametro, Calcular_Armadura_Minima, Dimensionar_Pilar_com_Esbelteza, Obter_Coeficiente_Beta, Calcular_Esbelteza, Calcular_Esbelteza_Limite, Calcular_Momento_Segunda_Ordem, Dimensionar_Secção_Rigoroso, Calcular_Armadura_Minima_Pilar, Otimizar_Armadura_Pilar, Dimensionar_Sapata_com_Dupla_Iteracao, Converter_Esforcos_Servico, Estimar_Dimensoes_Iniciais, Estimar_Peso_Proprio, Calcular_Tensoes_Solo, Arredondar_Superior, Calcular_Esforco_Puncoamento, Calcular_Resistencia_Puncoamento, Calcular_Armaduras_Flexao, Otimizar_Armadura_Sapata, Otimização, Armadura, Construtiva, Barras, Discretas, Metro, Obter_Area, ARREDONDAR_CIMA, Verificar_Espacamento, Encontrar_Menor_Area}
}
% --- FIM DOS ESTILOS DE CÓDIGO ---

% --- PACOTE PARA TABELAS ---
\usepackage{booktabs}

%Paginação
\usepackage{fancyhdr}
\pagestyle{fancy}
\fancyhf{}
\renewcommand{\headrulewidth}{0pt}
\fancyfoot[R]{\thepage}

%Criar os níveis de profundidade de subsubsubsection e de subsubsubsubsection
\setcounter{secnumdepth}{5}
\makeatletter
\newcommand\subsubsubsection{\@startsection{paragraph}{4}{\z@}{-2.5ex\@plus -1ex \@minus -.25ex}{1.25ex \@plus .25ex}{\normalfont\normalsize\bfseries}}
\newcommand\subsubsubsubsection{\@startsection{subparagraph}{5}{\z@}{-2.5ex\@plus -1ex \@minus -.25ex}{1.25ex \@plus .25ex}{\normalfont\normalsize\bfseries}}
\renewcommand\paragraph{\@startsection{paragraph}{4}{\z@}{-3.25ex\@plus -1ex \@minus -.2ex}{1.5ex \@plus .2ex}{\normalfont\normalsize\bfseries}}

\makeatother

%Indicar o nível de profundidade do índice geral
\setcounter{tocdepth}{5}

%Renomear o nome por defeito de secções padrão
\addto\captionsportuguese{\renewcommand{\contentsname}{ÍNDICE GERAL}}
\addto\captionsportuguese{\renewcommand{\listfigurename}{ÍNDICE DE FIGURAS}}
\addto\captionsportuguese{\renewcommand{\listtablename}{ÍNDICE DE TABELAS}}
\renewcommand{\lstlistingname}{Listagem}
\renewcommand{\lstlistlistingname}{ÍNDICE DE LISTAGENS}

%\addto\captionsportuguese{\renewcommand{\lstlistingname}{Listagem}} 


% As linhas seguintes foram movidas para o main.tex para os links funcionarem
% \addcontentsline{toc}{section}{ÍNDICE DE TABELAS}
% \addcontentsline{toc}{section}{ÍNDICE DE FIGURAS}
% \addcontentsline{toc}{section}{LISTA DE SIGLAS/ABREVIATURAS}
% \addcontentsline{toc}{section}{LISTA DE SÍMBOLOS}

% --- INÍCIO: PACOTE DE HYPERLINKS (DEVE SER DOS ÚLTIMOS A CARREGAR) ---
\usepackage[
    colorlinks=true,    % Usa cores em vez de caixas
    linkcolor=blue,     % Cor para links internos (Índice, figuras)
    citecolor=blue,     % Cor para links de bibliografia
    urlcolor=blue,      % Cor para URLs
    bookmarks=true,       % Cria marcadores no PDF
    bookmarksopen=true,   % Abre os marcadores por defeito
    pdfpagemode=UseOutlines % Mostra os marcadores em vez de miniaturas
]{hyperref}
% --- FIM DO PACOTE DE HYPERLINKS ---

%Colocar a numeração da secção à esquerda da numeração das equações, tabelas e figuras
\numberwithin{equation}{section}
\numberwithin{table}{section}
\numberwithin{figure}{section}

